
\documentclass{article}

% -----------------------------
% Packages
% -----------------------------
\usepackage{fancyhdr}
\usepackage{extramarks}
\usepackage{amsmath}
\usepackage{amsthm}
\usepackage{amsfonts}
\usepackage{tikz}
\usepackage[plain]{algorithm}
\usepackage{algpseudocode}
\usepackage{amssymb}

\usetikzlibrary{automata,positioning}

% -----------------------------
% Basic Document Settings
% -----------------------------
\topmargin=-0.45in
\evensidemargin=0in
\oddsidemargin=0in
\textwidth=6.5in
\textheight=9.0in
\headsep=0.25in

\linespread{1.1}

\pagestyle{fancy}
\lhead{\hmwkAuthorName}
\chead{\hmwkClass: \hmwkTitle}
\rhead{}
\lfoot{\lastxmark}
\cfoot{\thepage}

\renewcommand\headrulewidth{0.4pt}
\renewcommand\footrulewidth{0.4pt}

\setlength\parindent{0pt}

% -----------------------------
% Problem Section Helpers
% -----------------------------
\newcommand{\enterProblemHeader}[1]{
  \nobreak\extramarks{}{Problem #1 continued on next page\ldots}\nobreak{}
  \nobreak\extramarks{Problem #1 (continued)}{Problem #1 continued on next page\ldots}\nobreak{}
}

\newcommand{\exitProblemHeader}[1]{
  \nobreak\extramarks{Problem #1 (continued)}{Problem #1 continued on next page\ldots}\nobreak{}
  \nobreak\extramarks{Problem #1}{}\nobreak{}
}

\setcounter{secnumdepth}{0}
\newcounter{partCounter}
\newcounter{homeworkProblemCounter}
\setcounter{homeworkProblemCounter}{1}
\nobreak\extramarks{Problem \arabic{homeworkProblemCounter}}{}\nobreak{}

% -----------------------------
% Theorem env (optional)
% -----------------------------
\newtheorem*{theorem}{Theorem}

% -----------------------------
% Homework Problem Environment
%   Usage:
%   \begin{homeworkProblem}            % auto-number
%     % Your answer here
%   \end{homeworkProblem}
%
%   \begin{homeworkProblem}[1.2.3]     % label problem "1.2.3" (e.g. textbook numbering)
%     % Your answer here
%   \end{homeworkProblem}
% -----------------------------
\newenvironment{homeworkProblem}[1][]{
  \def\hwProblemLabel{#1}
  \ifx\hwProblemLabel\empty
    \edef\hwProblemLabel{\arabic{homeworkProblemCounter}}
    \stepcounter{homeworkProblemCounter}
  \fi
  \section{Problem \hwProblemLabel}
  \setcounter{partCounter}{1}
  \enterProblemHeader{\hwProblemLabel}
}{
  \exitProblemHeader{\hwProblemLabel}
}

% -----------------------------
% Homework Details (EDIT THESE)
% -----------------------------
\newcommand{\hmwkTitle}{Homework \#5}
\newcommand{\hmwkClass}{Math 3320}
\newcommand{\hmwkAuthorName}{\textbf{Joseph Quinn}}

% -----------------------------
% Title
% -----------------------------
\title{
  \vspace{2in}
  \textmd{\textbf{\hmwkClass:\ \hmwkTitle}}\\
  \vspace{3in}
}
\author{\hmwkAuthorName}
\date{}

% -----------------------------
% Convenience Macros (optional)
% -----------------------------
\renewcommand{\part}[1]{\textbf{\large Part \Alph{partCounter}}\stepcounter{partCounter}\\}

% Algorithms
\newcommand{\alg}[1]{\textsc{\bfseries \footnotesize #1}}

% Calculus
\newcommand{\deriv}[1]{\frac{\mathrm{d}}{\mathrm{d}x} (#1)}
\newcommand{\pderiv}[2]{\frac{\partial}{\partial #1} (#2)}
\newcommand{\dx}{\mathrm{d}x}

% Statistics
\newcommand{\E}{\mathrm{E}}
\newcommand{\Var}{\mathrm{Var}}
\newcommand{\Cov}{\mathrm{Cov}}
\newcommand{\Bias}{\mathrm{Bias}}

% Proof step (for aligned derivations)
\newcommand{\step}[2]{& #1 & & \text{#2} \\}

% Solution header (optional)
\newcommand{\solution}{\textbf{\large Solution}}

% -----------------------------
% Document
% -----------------------------
\begin{document}

\begin{homeworkProblem}[A10]

\end{homeworkProblem}
\[
	\begin{aligned}
		\max \quad & c_1 + 3c_3 + 5c_5                                                   \\
		\text{s.t.}\quad
		           & 4c_1 + 11c_2 - 16c_3 + 30c_4 + 14c_5 = 25  \quad             & (P1) \\
		           & c_1 + 3c_3 - 6c_4 + 7c_5 \leq 32                             & (P2) \\
		           & 5c_2 + 14c_3 - 9c_4 - c_5 \geq -8                            & (P3) \\
		           & c_1,c_2 \leq 0,\qquad c_3,c_4 \geq 0,\qquad c_5 \gtreqless 0
	\end{aligned}
\]

The goal is to find a linear combination of the primal constraints to create an upper bound on our $z$ objective function.


The first thing i will do is setup our inequality rectangle,

\begin{align*}
	4y_1c_1 + 11y_1c_2 - 16y_1c_3 + 30y_1c_4 + 14y_1c_5 & \le 25y_1 \\
	y_2c_1 + 3y_2c_3 - 6y_2c_4 + 7y_2c_5                & \le 32y_2 \\
	5y_3c_2 + 14y_3c_3 - 9y_3c_4 - y_3c_5               & \le -8y_3
\end{align*}

With,
\begin{align*}
	(4y_1 + y_2)c_1            & \ge 1c_1 \\
	(11y_1 + 5y_3)c_2          & \ge 0c_2 \\
	(-16y_1 + 3y_2 + 14y_3)c_3 & \ge 3c_3 \\
	(30y_1 - 6y_2 - 9y_3)c_4   & \ge 0c_4 \\
	(14y_1 + 7y_2 - y_3)c_5    & \ge 5c_5
\end{align*}

Our linear combination is our objective function and is written as,
\[
	w = y_{1}+y_{2}+y_{2} \ge 49
\]

We want to minimize $w$ as we want our bound to be as tight as possible.

The 49 comes from the sum of our RHS constants in our original constraints

Now to handle the inequalities and the variable constraints, we want to know what conditions on $y$ will make this work.
\\
For row 1, since (P1) has an equality relation, we know that row 1 works no matter what the negativty of $y_{1}$ is, meaning $y_{1} \gtrless$
\\
For row 2, since (P2) is $\le$ we know that only non negative values (or zero) for $y_{2}$ will make the row work as negative values would switch the sign. Meaning $y_{2} \ge 0$
\\
For row 3, since (P3) has a $\ge$ relation we know that in order for row 3 to work we need $y_{3}$ to be negative (or zero) in order for the sign to flip. Meaning $y_{3} \le 0$
\\
For col 1, since $c_{1} \le 0$ the inequality 	$(4y_1 + y_2)c_1 \ge 1c_1 $ can only be valid if we flip the sign to $\le$ as the $c_{1}$ value will flip it on division
\\
For col 2, since $c_{2} \le 0$ the inequality 	$	(11y_1 + 5y_3)c_2 \ge 0c_2 $ can only be valid if we flip the sign to $\le$ as the $c_{2}$ value will flip it on division
\\
For col 3, since $c_{3} \ge 0$ the inequality 	$		(-16y_1 + 3y_2 + 14y_3)c_3 \ge 3c_3 $ can only be valid if make the sign $\ge$ as the $c_{3}$ value will not flip it on division.
\\
For col 4, since $c_{4} \ge 0$ the inequality 	$		(30y_1 - 6y_2 - 9y_3)c_4 \ge 0c_4 $ can only be valid if make the sign $\ge$ as the $c_{4}$ value will not flip it on division.
\\
For col 5, since $c_{5} \ge 0$ the inequality 	$		(14y_1 + 7y_2 - y_3)c_5     \ge 5c_5$ can only be valid if make inequality an equality with the = sign, because $c_{5}$ is a free varaible and can be both positive, negative,
or zero. The only way the equation works is if we force an exact match.

Now we can then write the Dual LP problem as
\[
	\begin{aligned}
		\min \quad &
		w                      = 25y_1 + 32y_2 - 8y_3 \\
		\text{s.t.}\quad
		           & 4y_1 + y_2             \leq 1    \\
		           & 11y_1 + 5y_3           \leq 0    \\
		           & -16y_1 + 3y_2 + 14y_3  \geq 3    \\
		           & 30y_1 - 6y_2 - 9y_3    \geq 0    \\
		           & 14y_1 + 7y_2 - y_3     = 5       \\
		           & y_1 \gtreqless 0
		,y_2                    \geq 0
		y_3,                  \leq 0
	\end{aligned}
\]

\begin{homeworkProblem}[A11]

	\[
		\begin{align}
			\text{min}\quad
			w & = 5y_1 - 22y_3 + 8y_4                                      \\
			\text{s.t.}\quad
			  & y_1+9y_2+5y_3          = 12    &             & \text{(P1)} \\
			  & 5y_2+4y_3+2y_4         \le 18  &             & \text{(P2)} \\
			  & 2y_1-8y_2-5y_3+13y_4   \ge 0   &             & \text{(P3)} \\
			  & -6y_1+11y_4            \le -23 &             & \text{(P4)} \\
			  & 11y_1+11y_2+3y_3+8y_4          & \ge 15
			  &                                & \text{(P5)}
		\end{align}
	\]
	\[
		\begin{aligned}
			y_1 & \ge 0 &  & \text{(PS1)}, \\
			y_2 & \le 0 &  & \text{(PS2)}, \\
			y_3 & \le 0 &  & \text{(PS3)}, \\
			y_4 & \ge 0 &  & \text{(PS4)}.
		\end{aligned}
	\]

	Now having all the primal constraints and vairables named accoridng to the nmnenomic, we can determine the sign of each $x_m$ dual variable based on the primal constraint.

	Since the primal is a minimization problem,
	\[
		\begin{aligned}
			(P1):\ = \quad   & \Longrightarrow x_1 \text{ free}, \\
			(P2):\ \le \quad & \Longrightarrow x_2 \le 0,        \\
			(P3):\ \ge \quad & \Longrightarrow x_3 \ge 0,        \\
			(P4):\ \le \quad & \Longrightarrow x_4 \le 0,        \\
			(P5):\ \ge \quad & \Longrightarrow x_5 \ge 0.
		\end{aligned}
	\]
	Now we do the same thing but working from the other direction to determine each Dual constraints sign.

	Since the primal is a minimization problem,
	\[
		\begin{aligned}
			y_1\ge 0\ \text{(silly)}
			 & \Longrightarrow (D1):\ \le, \\
			y_2\le 0\ \text{(sensible)}
			 & \Longrightarrow (D2):\ \ge, \\
			y_3\le 0\ \text{(sensible)}
			 & \Longrightarrow (D3):\ \ge, \\
			y_4\ge 0\ \text{(silly)}
			 & \Longrightarrow (D4):\ \le.
		\end{aligned}
	\]
	The dual optimization function is created from the right hand constants of the primal making it,


	The dual objective is
	\[
		\boxed{
			\max z=12x_1+18x_2-23x_4+15x_5
		}
	\]
	Then the coefficient matrix is simply transposed making it,
	\[
		\begin{aligned}
			x_1+2x_3-6x_4+11x_5
			 & \le 5,   \\
			9x_1+5x_2-8x_3+11x_5
			 & \ge 0,   \\
			5x_1+4x_2-5x_3+3x_5
			 & \ge -22, \\
			2x_2+13x_3+11x_4+8x_5
			 & \le 8.
		\end{aligned}
	\]
	Now finally putting everything together we have the Dual LP,

	\[
		\begin{align}
			\max \quad z & = 12x_1 + 18x_2 - 23x_4 + 15x_5         \\
			\text{s.t.} \quad
			             & x_1 + 2x_3 - 6x_4 + 11x_5    \le 5      \\
			             & 9x_1 + 5x_2 - 8x_3 + 11x_5   \ge 0      \\
			             & 5x_1 + 4x_2 - 5x_3 + 3x_5    \ge -22    \\
			             & 2x_2 + 13x_3 + 11x_4 + 8x_5  \le 8      \\[4pt]
			             & x_1                          \gtreqless \\
			             & x_2                          \le 0      \\
			             & x_3                          \ge 0      \\
			             & x_4                          \le 0      \\
			             & x_5                          \ge 0
		\end{align}
	\]


\end{homeworkProblem}


\end{document}


